\begin{pastexam}[label={px:ncku-107-prog-05}]{成大資訊聯招 107 學年度 第 5 題}
（問答）

(10\%) What is the maximum numbers of elements in a heap of height $h$?
\end{pastexam}

\begin{pastsoln}
題目沒有定義高度。 這題的說法來自英文教科書（CLRS），那裡的高度數的是\textbf{邊數}：只有根的堆積
高度是 $0$。 資結書第~7.1~節的高度數的是\textbf{節點數}（層數），只有根的高度是 $1$。 下面兩種都寫，
作答時先寫一句「這裡的高度是邊數」。

答案：高度數邊數時最多 $2^{h+1} - 1$ 個元素；高度數節點數時最多 $2^{h} - 1$ 個。

\textbf{推導（高度數邊數）。} 二元堆積是完全二元樹（資結書第~8.1~節）。 根在深度 $0$，深度 $d$ 的
那一層最多 $2^d$ 個節點（每往下一層，每個節點最多兩個子節點，個數最多加倍）。 高度 $h$ 表示最深的
節點深度是 $h$，所以節點都在深度 $0, 1, \ldots, h$，總數最多
\[
\sum_{d=0}^{h} 2^d = 2^{h+1} - 1 .
\]
每一層都排滿的完全二元樹（滿二元樹）正好有這麼多個節點，所以上界達得到，最大值就是 $2^{h+1} - 1$。

\textbf{高度數節點數時}，高度 $h$ 的樹只有深度 $0..h-1$ 這 $h$ 層，同樣的算法得到
$\sum_{d=0}^{h-1} 2^d = 2^h - 1$。 兩個答案只差在 $h$ 的意思，換算是「邊數高度 $=$ 節點數高度 $- 1$」。

例：$h = 2$（邊數）時，深度 $0, 1, 2$ 三層排滿是 $1 + 2 + 4 = 7 = 2^3 - 1$ 個。
\end{pastsoln}
